\selectlanguage{ngerman}%
\frenchspacing%
\begin{ledgroupsized}[r]{120mm}%
\footnotesize%
\pstart%
\noindent\textbf{Überlieferung:}
\pend%
\end{ledgroupsized}%
\begin{ledgroupsized}[r]{114mm}%
\footnotesize%
\pstart%
\parindent -6mm%
\makebox[6mm][l]{\textit{L}}%
Aufzeichnung:
LH~XXXV~11,\;14 Bl.\;31. 
Ein Zettel (etwa 6,5\;x\;8,5~cm).
Eineinhalb Seiten;
Textfolge: Bl.\;31\,v\textsuperscript{o}, 31\,r\textsuperscript{o}, durch Kustos (\textit{verte}) festgelegt.
Auf Bl.\;31\,r\textsuperscript{o}, senkrecht zur Schreibrichtung, ist das Überbleibsel eines nicht weiter ermittelten, bereits vorhandenen Textes von Leibnizens Hand (wohl das Fragment eines Briefkonzeptes) anzutreffen:
\textlangle...\textrangle\textit{ster,} %
\lbrack/\rbrack\
\textlangle...\textrangle\textit{er bewusten erkundigung}
\lbrack/\rbrack\
\textlangle...\textrangle\textit{ncke mich dienstlich}
\lbrack/\rbrack\
\textlangle...\textrangle\
\textit{Ludolphi, und dem Hr.}%
\indexN{\textls{Ludolph} (Ludolphus), Hiob 1624\textendash1704}
\lbrack/\rbrack\
\textlangle...\textrangle\
\textit{bey gelegenheit meines}
\lbrack/\rbrack\
\textlangle...\textrangle\
\textit{und werden Sie bald}
\textlangle...\textrangle%
\pend%
\end{ledgroupsized}%
\vspace{5mm}%
\begin{ledgroup}%
\footnotesize%
\pstart%
\noindent%
\textbf{Datierungsgründe:}
Im Texträger der titellosen und undatierten Aufzeichnung N.\;\ref{58946}, die einer Klassifizierung der Stoßfälle gewidmet ist, liegt kein Wasserzeichen vor.
Bei der im deutschen Textfragment erwähnten Person handelt es sich sehr wahrscheinlich um Hiob Ludolph,\indexN{\textls{Ludolph} (Ludolphus), Hiob 1624\textendash1704}
mit dem Leibniz jahrelang korrespondierte
(die edierte Korrespondenz erstreckt sich nämlich von Dezember 1687 bis Oktober 1703; Hiob Ludolph starb 1704).
Das betreffende Fragment entstammt aber unwahrscheinlich einem an Ludoph selbst gerichteten Brief, der ohnehin in Lateinischem verfasst worden wäre.
Da der Textträger wohl kein italienisches Papier ist, wurde der Text mit großer Wahrscheinlichkeit nicht zu der Zeit verfasst, als Leibniz sich in Italien befand (von März 1689 bis März 1690).
\pend%
\pstart%
Im edierten Briefwechsel lässt sich kein Brief auffindig machen, der mit unmittelbarer Evidenz dem Fragment in N.\;\ref{58946} entspricht.
Am meisten könnte das deutschsprachige Konzept \textit{LSB} I,\;14 N.\;319, S.\;542\,f.\cite{01424} (Leibniz für Nicolaus Förster, \lbrack 9.\;Oktober\rbrack\ 1697) in Frage kommen:\indexN{\textls{Förster}, Nicolaus 1657\textendash1732}
Das \glqq Memorial\grqq\ sollte als Grundlage für einen Brief an För\textit{ster} (ebd., S.\;542.5) dienen;
Otto M\textit{encke} wird am Anfang auch als Adressat erwähnt (ebd., Z.\;9);\indexN{\textls{Mencke} (Menken, Menkenius, Menque), Otto 1644\textendash1707}
kurz danach wird auch Hiob Ludolph im Zusammenhang mit anderen Gelehrten erwähnt (ebd., Z.\;13).
Ein noch schwächerer Kandidat ware das deutschsprachige Briefkonzept \textit{LSB} I,\;10 N.\;61, S.\;74\,f.\cite{01423} (Leibniz an Ludolf Hugo, 4./14.\;Oktober 1694).\indexN{\textls{Hugo}, Ludolf 1632\textendash1704}
Ansonsten ist festzustellen, dass in den 1690er Jahren die Erwähnungen von Hiob Ludolf in Leibnizens Briefwechsel deutlich dichter sind als zuvor (die meisten Erwähnungen finden sich in Bänden von \textit{LSB},\;I).\indexN{\textls{Ludolph} (Ludolphus), Hiob 1624\textendash1704}
Das Brieffragment ist somit unwahrscheinlich vor 1690 bzw. vor Leibnizens Abreise aus Italien verfasst worden.
Davon rührt die vorgeschlagene Datierung her.
Die Aufzeichnung N.\;\ref{58946} könnte demnach im Zusammenhang mit der unvollendeten Stoßlehre von Leibnizens \textit{Dynamica}\cite{01354} entstanden sein (siehe hierüber die einleitende Erläuterung zum Textkomplex N.\;\ref{impulsuseccentricus}, S.\;\pageref{impulsecc_intro}\,ff.).
\pend%
\pstart%
Als datierungsrelevant könnten sich ferner folgende Umstände erweisen: \ 
(1)~%
Nach der in N.\;\ref{58946} vorgeschlagenen Klassifizierung der Stoßfälle wäre es unangemessen, von einem \textit{direkten} exzentrischen Stoß zu reden, wie dies in N.\;\ref{60142+60297}, \refpassage{60142+60297_anfang-1}{60142+60297_anfang-1} geschieht.
Dies könnte zeigen, dass N.\;\ref{58946} und N.\;\ref{60142+60297} unwahrscheinlich in unmittelbarer Reihenfolge entstanden. \ 
(2)~%
Der Schlussteil von N.\;\ref{58946} (\refpassage{LH_35_11_14_031_031v_tang-1}{LH_35_11_14_031_031v_tang-2}) könnte an die Frage zur \textit{obliquitas concursus} anknüpfen, die in knapper Weise in der Aufzeichnung N.\;\ref{58950} (\refpassage{LH_35_11_14_035_035v_tang-1}{LH_35_11_14_035_035v_tang-2}) gestellt wird.
Diese letztere ist \textendash\ unter Vorbehalt \textendash\ auf die Zeit nach dem 10.\;November 1690 datiert.
(3)~%
Inhaltliche Verwandschaft (hinsichtlich des \textit{punctum contactus} bzw. \textit{planum contactus}) könnte auch mit dem Text N.\;\ref{58293} (\textit{Ex conservatione potentiae, tam absolutae quam respectivae, possunt derivari phaenomena concursus trium corporum}) im vorliegenden Band bestehen, der wohl nicht vor dem 8.\;Mai 1693 entstanden sein dürfte.
\pend 
\end{ledgroup}
\selectlanguage{latin}
\frenchspacing
\newpage%
\pstart%
\normalsize%
\noindent%
\lbrack%
\href{\Scan{LH_35_11_14_031v}}%
{31\,v\textsuperscript{o}}%
\rbrack\ %
\edtext{%
Omnis%
\textls{ concursus }%
\indexS{concursus corporum}%
est vel simplex%
\indexS{concursus simplex}
vel compositus.%
\indexS{concursus compositus}%
}{%
\lemma{Omnis}\Bfootnote{%
\textit{(1)}~\textls{incursus}
\textit{(2)}~\textls{concursus} est
\textit{(a)}~vel directus, vel indirectus.%
\indexS{incursus directus}\indexS{incursus indirectus}
\textit{(b)}~vel simplex vel compositus.%
~\textit{L}}}%
\edtext{%
\textls{ Simplex }%
\indexS{concursus simplex}%
si unus tantum sit contactus,%
\indexS{contactus unus}
quod fit
sive unum tantum sit punctum contactus,%
\indexS{punctum contactus}
sive una linea,%
\indexS{linea contactus}
aut una superficies.%
\indexS{superficies contactus}%
\textls{ Compositus }%
\indexS{concursus compositus}%
}{%
\lemma{\textls{Simplex}}\Bfootnote{%
\textit{(1)}~cum
\textit{(a)}~\textlangle\textendash\textrangle\
\textit{(b)}~una
\textit{(aa)}~\textlangle\textendash\textrangle\
\textit{(bb)}~est tantum
\textit{(aaa)}~\textlangle\textendash\textrangle\ %
\textit{(bbb)}~ictus plaga.\indexS{plaga ictus} Ut si
\textit{(2)}~si unum sit corpus incurrens
\textit{(3)}~si duo
\textit{(4)}~si unus tantum sit contactus,
\textit{(a)}~compositus
\textit{(b)}~quod fit sive unum
\bars{sit \textit{streicht Hrsg.}} tantum sit %
\lbrack...\rbrack\ superficies. \textls{Compositus}%
~\textit{L}}}%
sive plura sint corpora quam duo,%
\indexS{corpora concurrentia}
sive duorum corporum plures sint contactus discontinuati.%
\indexS{contactus plures discontinuati}
\pend%
\pstart%
Contactus%
\textls{ simplex }%
est vel%
\textls{ directus }%
\indexS{contactus simplex directus}%
vel%
\textls{ indirectus.}%
\indexS{contactus simplex indirectus}
Contactus%
\textls{ directus }%
\indexS{contactus simplex directus}%
est cum linea perpendicularis ad contactum%
\indexS{linea perpendicularis ad contactum}
transit per centra gravitatis duorum corporum%
\indexS{centra gravitatis corporum concurrentium}
\edtext{%
concurrentium;%
\indexS{corpora concurrentia}
potest etiam%
}{%
\lemma{concurrentium;}\Bfootnote{%
\textit{(1)}~et
\textit{(2)}~potest etiam%
~\textit{L}}}
definiri ab
\edtext{%
effectu,%
\indexS{effectus}
cum linea directionis%
\indexS{linea directionis mutata concursu}
punctorum in corporibus%
\indexS{corpora concurrentia}
quae erat ante concursum,
concursu%
\indexS{concursus corporum}
non mutatur.%
}{%
\lemma{effectu,}\Bfootnote{%
\textit{(1)}~cum contactus non mutat, \bars{cum \textit{erg.}} linea
\textit{(2)}~cum
\textit{(3)}~cum linea %
\lbrack...\rbrack\ in corporibus
\textit{(a)}~contactu non mutatur
\textit{(b)}~quae erat %
\lbrack...\rbrack\ non mutatur.%
~\textit{L}}}
Quanquam
\edtext{enim}{%
\lemma{enim}\Bfootnote{%
\textit{erg.\;L}}}
hoc contingere possit
etiam in concursu indirecto%
\indexS{concursus indirectus}
sed composito,%
\indexS{concursus compositus}
tamen inter simplices%
\indexS{concursus simplex}
soli competit
\edtext{%
directo.%
\indexS{concursus directus}
Directi igitur requisita sunt,%
\indexS{requisitum}
ut sit simul et perpendicularis%
\indexS{concursus perpendicularis}
et centricus.%
\indexS{concursus centricus}%
\textls{ Indirectus }fit,
sive non sit perpendicularis%
\indexS{concursus perpendicularis}
sed obliquus,%
\indexS{concursus obliquus}
sive non sit centricus%
\indexS{concursus centricus}
sed eccentricus.%
\indexS{concursus eccentricus}
Itaque est vel obliquus tantum,%
\indexS{concursus obliquus}
vel eccentricus tantum,%
\indexS{concursus eccentricus}
vel simul et eccentricus%
\indexS{concursus eccentricus}
et obliquus.%
\indexS{concursus obliquus}%
}{%
\lemma{directo.}\Bfootnote{%
\textit{(1)}~Indirectus est
\textit{(a)}~vel
\textit{(b)}~vel
\textit{(aa)}~eccen
\textit{(bb)}~obliquus centricus
\textit{(aaa)}~vel ec
\textit{(bbb)}~sive
\textit{(2)}~Directi igitur %
\lbrack...\rbrack\ centricus. \textls{Indirectus}
\textit{(a)}~vel est non perpendicularis
\textit{(b)}~fit, sive %
\lbrack...\rbrack\ obliquus tantum,
\textit{(aa)}~vel centricus %
\textit{(bb)}~vel eccentricus %
\lbrack...\rbrack\ et obliquus.%
~\textit{L}}}
\lbrack%
\href{\Scan{LH_35_11_14_031r}}%
{31\,r\textsuperscript{o}}%
\rbrack%
\pend%
\pstart%
\edlabel{LH_35_11_14_031_031v_tang-1}%
Re attentius considerata,%
\indexS{res considerata}
melius definiemus concursum simplicem,%
\indexS{concursus simplex}
ut sit is
in quo una tantum est linea perpendicularis ad contactum,%
\indexS{linea perpendicularis ad contactum}
\edtext{%
idque fiet%
}{\lemma{idque}\Bfootnote{%
\textit{(1)}~erit
\textit{(2)}~fiet%
~\textit{L}}}
non in puncto tantum,
sed et in tali puncto %
in quo non nisi una dici potest superficies%
\indexS{superficies plana tangens}
\edtext{plana}{%
\lemma{plana}\Bfootnote{%
\textit{erg.\;L}}}
tangens%
\lbrack:\rbrack\
nam interdum corpus in eodem puncto%
\indexS{corpus tactum a pluribus planis in puncto}
a pluribus planis%
\indexS{planum tangens}
tangi potest.%
\edlabel{LH_35_11_14_031_031v_tang-2}
Ex.\;gr. cycloidalis linea%
\indexS{linea cycloidalis}
in 
\edtext{%
extremo puncto%
\indexS{punctum extremum}
ubi%
}{%
\lemma{extremo}%
\Bfootnote{%
\textit{(1)}~ubi
\textit{(2)}~puncto ubi%
~\textit{L}%
}}
basin%
\indexS{basis}
attingit.
Et ita
quia unica hoc modo est linea perpendicularis,%
\indexS{linea perpendicularis ad contactum}
possumus deinde considerare
utrum ea cum linea ictus%
\indexS{linea ictus}
coincidat
et utrum transeat
\edtext{%
per centra%
}{\lemma{per}\Bfootnote{%
\textit{(1)}~centrum
\textit{(2)}~centra%
~\textit{L}}}
gravitatis.%
\indexS{centra gravitatis corporum concurrentium}
\pend%
\selectlanguage{latin}%
\frenchspacing%
\normalsize%
